\documentclass[12pt]{article}

\usepackage[a4paper, margin=1in]{geometry}
\usepackage{graphicx, parskip, setspace, ragged2e, enumitem}
\usepackage{amsmath, amssymb, amsthm, tikz, tikz-3dplot, chemfig, physics, cancel, overarrows}
\usepackage[english]{babel}
\setstretch{1.5}

\usepackage{accents}

\renewcommand{\vec}[1]{\undertilde{#1}}

\begin{document}
Suppose a sphere $\mathcal{S}$ of radius $1$ centred at point $P$ with position vector \[\vec{p}=\begin{bmatrix}
		10 \\
		0  \\
		2
	\end{bmatrix},\]
and define the line
\[\ell=\begin{bmatrix}
		0 \\
		0 \\
		2
	\end{bmatrix} + k\begin{bmatrix}
		1 \\
		Y \\
		Z
	\end{bmatrix},\]
  where $k\in \mathbb{R}$ and $Y, Z \sim \mathcal{U}[-1,1]$.
\vskip 8pt
\hrule
\begin{enumerate}
	\item[(i)] By considering an appropriate perpendicular-distance expression, deduce that for $\ell$ to intersect $\mathcal{S}$,
	      \[\dfrac{10\sqrt{Y^2+Z^2}}{\sqrt{1+Y^2+Z^2}} \leq 1,\]
	      and further show that
	      \[Y^2 + Z^2 \leq \dfrac{1}{99}.\]
	\item[(ii)] Explain why this result is meaningful.
	      \textbf{You may assume the fact that for $X, Y \sim \mathcal{U}[-n,n], n \in \mathbb{R}^+$, the point $P(X, Y)$ is uniformly distributed in the resulting square without proof.}
	\item[(iii)] Hence, by considering the ratio of favourable outcomes to total outcomes, deduce the probability that the line $\ell$ intersects $\mathcal{S}$.
\end{enumerate}
\end{document}
